% ---------------------------------------------------------------------
%  §6. Reflexión metodológica
% ---------------------------------------------------------------------
\section{Reflexión metodológica}

\textbf{Uso de IA.} La IA se utilizó como asistente para sugerir transformaciones y gráficos, revisar la
redacción de las interpretaciones y redactar este informe en \LaTeX{} a partir del notebook ejecutado, con
una instrucción de rol de analista de datos que exigía usar exactamente sus números, gráficos y conclusiones.
La interpretación y las decisiones de análisis fueron responsabilidad del grupo.
\textbf{Revisiones.} Cada integrante contrastó las afirmaciones del informe con el notebook y los datos
originales, y el grupo consolidó los comentarios. [COMPLETAR: comentarios de la validación externa y cómo se
incorporaron.]
\textbf{Aprendizajes.} Auditar la calidad de los datos es indispensable: la verificación estándar no detectaba
filas repetidas, y solo al notar que el identificador máximo (1.000) era menor que la cantidad de filas
(1.743) se descubrieron 705 estudiantes cargados dos veces. También se comprobó que una diferencia entre
grupos puede deberse a una tercera variable —las horas de uso—, lo que cambió la lectura de las diferencias
entre plataformas.
